% Chương 5: Hướng phát triển

\chapter{Hướng phát triển}

Dựa trên kết quả và hạn chế hiện tại, đồ án có thể được mở rộng và cải tiến theo các hướng sau:

\textbf{Xử lý streaming thời gian thực:} Chuyển từ batch processing sang streaming architecture sử dụng Spark Structured Streaming hoặc Apache Flink. Điều này cho phép phát hiện xu hướng với độ trễ chỉ vài phút thay vì bốn giờ, đặc biệt quan trọng với các xu hướng viral nhanh. Cần thiết kế lại pipeline để xử lý incremental updates và maintain state cho clustering.

\textbf{Mở rộng nguồn dữ liệu:} Tích hợp thêm các nền tảng mạng xã hội lớn như Facebook, Instagram, Twitter/X, và các nền tảng Việt Nam như Zalo, Threads. Mỗi nền tảng có đặc thù riêng về API và loại nội dung, yêu cầu adapter pattern linh hoạt. Ngoài ra, có thể mở rộng sang các nguồn dữ liệu khác như podcast, forum, e-commerce reviews.

\textbf{Tối ưu chi phí:} Giảm phụ thuộc vào API trả phí bằng cách: (1) self-host embedding models như sentence-transformers hoặc Gemma; (2) sử dụng mô hình ngôn ngữ lớn mã nguồn mở như Llama 3 hoặc Qwen cho analysis; (3) implement caching thông minh để tránh tính toán lại. Trade-off là cần đầu tư infrastructure (GPU) nhưng chi phí dài hạn thấp hơn.

\textbf{Cải thiện quality filtering:} Áp dụng các kỹ thuật machine learning để phát hiện spam và bot phức tạp hơn. Huấn luyện classifier dựa trên features như writing patterns, posting frequency, account age. Tích hợp sentiment analysis để lọc nội dung tiêu cực hoặc toxic trước khi phân cụm.

\textbf{Dự đoán xu hướng:} Xây dựng mô hình time-series forecasting để dự đoán xu hướng sắp bùng nổ dựa trên tốc độ tăng trưởng engagement, pattern lan truyền, và historical data. Sử dụng kỹ thuật như LSTM, Prophet, hoặc Transformer-based models. Điều này cho phép doanh nghiệp chuẩn bị chiến lược proactive thay vì reactive.

\textbf{Personalization:} Xây dựng hệ thống recommendation cá nhân hóa, gợi ý xu hướng phù hợp với từng user dựa trên industry, interests, past interactions. Áp dụng collaborative filtering hoặc content-based filtering kết hợp với embedding similarity.

\textbf{Hỗ trợ đa ngôn ngữ:} Mở rộng sang các ngôn ngữ khác ngoài tiếng Việt và tiếng Anh. Sử dụng multilingual embeddings như mBERT hoặc XLM-RoBERTa để xử lý nhiều ngôn ngữ trong cùng một không gian vector, cho phép phát hiện xu hướng cross-language.

\textbf{Visual content analysis:} Tích hợp xử lý hình ảnh và video để phân tích visual trends. Sử dụng CLIP hoặc các mô hình multimodal để extract embeddings từ thumbnail, video frames, kết hợp với text embeddings để có phân tích toàn diện hơn.

\textbf{Monitoring và alerting:} Xây dựng hệ thống monitoring chi tiết với metrics (processing time, API costs, cluster quality), dashboards cho admin, và alerting khi phát hiện anomalies hoặc xu hướng rủi ro cao. Tích hợp với Prometheus, Grafana hoặc các công cụ APM.

\textbf{API và integration:} Cung cấp RESTful API hoặc GraphQL endpoint cho phép các ứng dụng bên ngoài tích hợp dữ liệu xu hướng. Xây dựng webhook để notify real-time khi phát hiện xu hướng mới. Điều này mở ra khả năng commercialize hệ thống như một trend detection service.
